%% ---------------------------------------------------------------------------
\section{Applications to Four and Five States}\label{sec:applications}

We now combine the necessary conditions with satisfiability solvers. All
statements of this section are computer-assisted. Each of them is either a
finite simulation or enumeration, or the unsatisfiability of a propositional
formula that is certified by a proof checked independently of the solver; the
kind of evidence is stated for each result and summarized in
\cref{tab:results}.

\subsection{Method}\label{sec:method}

\paragraph{Encodings}
The formula $\MT_k(\mathcal N;M,N')$ has one-hot variables for the entries of a
$k$-state rule table, for the cells of $C_\infty$ with $t+i\le M$ and for the
cells of $R_n$ with $n\in\mathcal N$. Its clauses state the local rule, that no
half-line cell is $\Fq$, that the line of length $n$ is synchronized at time
$2n-2$ for every $n\in\mathcal N$, the return chains of \cref{lem:return} for
$n,n+1\in\mathcal N$, and $\PUMP(n,n')$ for $4\le n<n'\le N'$; the formula
$\MT_k(\mathcal N;M)$ has no pumping clauses. The formula $\G_k(T,N_1,C)$ is
encoded in the same way, with $M=2N_1$. Its regions are closed under
predecessors, so the encoding is equivalent to \cref{def:G}.

\paragraph{Certificates}
Every unsatisfiability claim is certified by a proof in the LRAT
format~\cite{CruzFilipeEtAl17}, which is produced by the solver
CaDiCaL~\cite{CaDiCaL2} or, for one formula, obtained from a DRAT proof of the
solver Kissat~\cite{Kissat20} with \textsf{drat-trim}~\cite{WetzlerHeuleHunt14},
and verified by the independent checker \texttt{lrat-check}. The formula
generators, the enumerations and the simulations are trusted programs, and
every model found by a solver was re-simulated by a checker that shares no code
with the generators.

\paragraph{Validation}
As tests of the encodings, Mazoyer's solution, as extracted from Duprat's
formal proof of its correctness~\cite{Duprat97}, satisfies
$\MT_6(\{3,\dots,12\};78,40)$ and, with its table fixed, (G2) and (G3) for $T=8$
and all $n\ge15$. With the table of the rule $\delta_{14}$ of \cref{sec:five}
fixed, the formulas are satisfiable for the lines of lengths $2,\dots,14$ and
unsatisfiable with the line of length $15$, or with the early rows ($T=8$, but
not $T=6$) of the lines of lengths $15,\dots,30$, as simulation predicts.

\paragraph{Reproducibility}
The supplementary material contains the programs, the rule tables, the lists of
half-line germs and the logs, with two scripts that rerun every computation:
one for the simulations, enumerations and satisfiable formulas (about $7$
minutes), and one that regenerates and checks every LRAT proof (about $30$
minutes). The solver times, proof sizes and SHA-256 digests of all formulas and
proofs are listed there; see the data availability statement at the end.

\subsection{Four States}\label{sec:four}

\begin{theorem}\label{thm:four}
No four-state rule satisfying \eqref{eq:quiescence}
\begin{enumerate}[(i)]
\item synchronizes all lines of lengths $2,\dots,9$ in minimal time;
\item satisfies $\G_4(6,10)$: the lines of lengths $2,\dots,7$ complete and,
for the lines of lengths $8$, $9$ and $10$, the first seven rows and the return
chains;
\item satisfies $\G_4(4,10,10)$: the lines of lengths $2,\dots,5$ complete and,
for the lines of lengths $6,\dots,10$, the first five rows, the return chains
and the right ends.
\end{enumerate}
In particular, no four-state minimal-time solution exists.
\end{theorem}

\begin{proof}
(i) Such a rule yields a model of $\MT_4(\{2,\dots,9\};16)$: a half-line cell
with $t+i\le16$ is quiescent if $i\ge10$ (\cref{lem:lightcone}) and equals a
cell of the line of length $9$ at a time $t<16$ if $i\le9$ (\cref{lem:agree}),
so it is in $W$; and the return chains hold by \cref{lem:return} for $2\le
n\le8$. This formula has $753$ variables and $17{,}605$ clauses. (ii) The
encoding of $\G_4(6,10)$ has $919$ variables and $22{,}960$ clauses. (iii) The
encoding of $\G_4(4,10,10)$ has $874$ variables and $20{,}461$ clauses. CaDiCaL
refutes each of the three formulas with an LRAT proof, and \texttt{lrat-check}
verifies each proof. The last claim follows from (i), and also from (ii) or
(iii) with \cref{lem:G}.
\end{proof}

Part (i) certifies the bound observed by Yun\`es~\cite{Yunes93}; the earlier
proofs~\cite{Balzer67,Yunes93,Sanders94} rest on the correctness of search
programs. Parts (ii) and (iii) show how little information excludes four
states: in (iii), only the lines of lengths up to $5$ are complete. These
bounds are sharp in the following sense.

\begin{proposition}[computational]\label{prop:sharp}
\leavevmode
\begin{enumerate}[(a)]
\item Exactly $27$ partial four-state rules (rules restricted to the
neighbourhoods they use) synchronize all lines of lengths $2,\dots,8$ in
minimal time, and no partial four-state rule synchronizes all lines of lengths
$2,\dots,9$.
\item The formulas $\G_4(5,30)$, $\G_4(6,9)$, $\G_4(4,30,9)$ and
$\G_4(3,30,12)$, the formula $\G_4(6,10)$ without (G3), and the formula that
consists of (G1) for $T=4$ and (G5) for $6\le n\le12$ are satisfiable.
\end{enumerate}
\end{proposition}

\begin{proof}
(a) is an exhaustive enumeration. Yun\`es~\cite{Yunes93} found the same number
under the extra condition $\delta(\ast,\Lq,\Lq)=\Lq$, which is vacuous here
because none of the $27$ partial rules uses the neighbourhood $(\ast,\Lq,\Lq)$.
(b) The solver finds a model of each formula, and each model is re-simulated.
\end{proof}

So in part (ii) of \cref{thm:four} neither $T$ nor $N_1$ can be lowered and
(G3) cannot be dropped; in part (iii) neither $T$ nor $C$ can be lowered, even
with more early rows, and neither the early rows nor the right ends suffice on
their own.

\subsection{Five States: Rules for the Lengths Up to Fourteen}\label{sec:five}

The following five-state rules were found with satisfiability solvers. Their
tables are the files \texttt{uniformA\_2-12.txt} (the rule $\delta_{12}$),
\texttt{delta13.txt} and \texttt{delta14*.txt} in \texttt{results/} of the
supplementary material. \Cref{tab:delta14} shows $\delta_{14}$.

\begin{table}[t]
\centering\small
\caption{The five-state rule $\delta_{14}$. The entry in row $(x,y)$ and
column $z$ is $\delta_{14}(x,y,z)$; the neighbourhoods $(\ast,y,\ast)$ are not
used (``--'').}
\label{tab:delta14}
\setlength{\tabcolsep}{4.2pt}
\begin{tabular}{@{}cc|ccccc@{\hspace{2.2em}}cc|ccccc@{}}
\toprule
$x$ & $y$ & $z=\Lq$ & $\Gq$ & $\Aq$ & $\Bq$ & $\ast$ &
$x$ & $y$ & $z=\Lq$ & $\Gq$ & $\Aq$ & $\Bq$ & $\ast$\\
\midrule
%% generated by figures.py from results/delta14.txt; do not edit
$\ast$ & $\mathsf L$ & $\mathsf F$ & $\mathsf F$ & $\mathsf L$ & $\mathsf G$ & -- & $\ast$ & $\mathsf A$ & $\mathsf G$ & $\mathsf G$ & $\mathsf L$ & $\mathsf G$ & --\\
$\mathsf L$ & $\mathsf L$ & $\mathsf L$ & $\mathsf G$ & $\mathsf G$ & $\mathsf B$ & $\mathsf L$ & $\mathsf L$ & $\mathsf A$ & $\mathsf B$ & $\mathsf G$ & $\mathsf B$ & $\mathsf A$ & $\mathsf A$\\
$\mathsf G$ & $\mathsf L$ & $\mathsf A$ & $\mathsf B$ & $\mathsf A$ & $\mathsf B$ & $\mathsf G$ & $\mathsf G$ & $\mathsf A$ & $\mathsf A$ & $\mathsf G$ & $\mathsf L$ & $\mathsf A$ & $\mathsf B$\\
$\mathsf A$ & $\mathsf L$ & $\mathsf B$ & $\mathsf G$ & $\mathsf G$ & $\mathsf G$ & $\mathsf G$ & $\mathsf A$ & $\mathsf A$ & $\mathsf L$ & $\mathsf A$ & $\mathsf G$ & $\mathsf A$ & $\mathsf B$\\
$\mathsf B$ & $\mathsf L$ & $\mathsf B$ & $\mathsf A$ & $\mathsf A$ & $\mathsf L$ & $\mathsf A$ & $\mathsf B$ & $\mathsf A$ & $\mathsf A$ & $\mathsf A$ & $\mathsf G$ & $\mathsf A$ & $\mathsf L$\\
\midrule
$\ast$ & $\mathsf G$ & $\mathsf L$ & $\mathsf F$ & $\mathsf L$ & $\mathsf L$ & -- & $\ast$ & $\mathsf B$ & $\mathsf G$ & $\mathsf G$ & $\mathsf G$ & $\mathsf G$ & --\\
$\mathsf L$ & $\mathsf G$ & $\mathsf A$ & $\mathsf F$ & $\mathsf B$ & $\mathsf A$ & $\mathsf F$ & $\mathsf L$ & $\mathsf B$ & $\mathsf B$ & $\mathsf G$ & $\mathsf A$ & $\mathsf L$ & $\mathsf G$\\
$\mathsf G$ & $\mathsf G$ & $\mathsf B$ & $\mathsf F$ & $\mathsf G$ & $\mathsf F$ & $\mathsf B$ & $\mathsf G$ & $\mathsf B$ & $\mathsf B$ & $\mathsf A$ & $\mathsf B$ & $\mathsf B$ & $\mathsf F$\\
$\mathsf A$ & $\mathsf G$ & $\mathsf A$ & $\mathsf G$ & $\mathsf G$ & $\mathsf A$ & $\mathsf B$ & $\mathsf A$ & $\mathsf B$ & $\mathsf B$ & $\mathsf B$ & $\mathsf G$ & $\mathsf B$ & $\mathsf G$\\
$\mathsf B$ & $\mathsf G$ & $\mathsf A$ & $\mathsf G$ & $\mathsf A$ & $\mathsf L$ & $\mathsf B$ & $\mathsf B$ & $\mathsf B$ & $\mathsf L$ & $\mathsf A$ & $\mathsf A$ & $\mathsf L$ & $\mathsf A$\\
\bottomrule

\end{tabular}
\end{table}

\begin{proposition}\label{prop:nearmiss}
Each of the five-state rules $\delta_{14}$, $\delta_{14}^{\mathrm b}$,
$\delta_{14}^{\mathrm c}$, $\delta_{14}^{\mathrm d}$ and $\delta_{14}^{\mathrm
e}$ synchronizes the line of length $n$ at time $2n-2$ for $2\le n\le14$, and
for no $n\ge15$. More precisely, for every $n\ge15$, some cell of the line of
length $n$ is $\Fq$ at a relative time $\tau\le\min(T,n-2)$, where
$T=10,14,22,12,14$, respectively.
\end{proposition}

\begin{proof}
The lines of lengths $2,\dots,60$ are simulated. For $15\le n\le60$ some cell
is $\Fq$ at a relative time $\tau\le n-2$, which proves the claim for $15\le
n<T+3$. The half-line has no cell in state $\Fq$ at the times $t<3000$, so by
\cref{lem:agree} a cell of a line of length $n\le1500$ that is $\Fq$ at a time
$t\le 2n-2$ lies in $R_n$. Simulating the half-line gives $t_0$ and $p$ that
satisfy the hypothesis of \cref{lem:certificate} for $D=2T+4\ge2T+1$, namely
$(t_0,p)=(51,8)$, $(66,3)$, $(257,192)$, $(95,48)$ and $(56,2)$. For every $n$
with $\max(15,T+3)\le n<t_0+3+p\le452$, direct simulation shows that the line
of length $n$, or of a length $n'$ in this range with $J_{n'}=J_n$ on
$[0,T+1]$, has a cell in state $\Fq$ at a relative time $\tau\le T$, hence in
$E(T)$. (One length is simulated for each class of a partition that is finer
than equality of $J_n$ on $[0,T+1]$; together with the lengths below $T+3$,
these are $35$, $39$, $413$, $116$ and $28$ lengths, the largest being $49$,
$53$, $427$, $130$ and $42$.) By \cref{lem:early,cor:reduce} with
$N_0=\max(15,T+3)$, every $R_n$ with $n\ge N_0$ has a cell of $E(T)$ in state
$\Fq$, and there $\tau\le T\le n-3$.
\end{proof}

\begin{figure}[t]
\centering
%% Figure 5: delta14 on the lines of lengths 14 and 15, and its half-line (cells from figures.py).
\begin{tikzpicture}
  % the line of length 14, up to its firing time 26
  \begin{scope}[x=1.55mm,y=1.55mm]
    %% generated by figures.py; do not edit
\fill[stG] (0,0) rectangle (1,-1);
\fill[stA] (1,-1) rectangle (2,-2);
\fill[stB] (1,-2) rectangle (3,-3);
\fill[stG] (0,-3) rectangle (1,-4);
\fill[stB] (3,-3) rectangle (4,-4);
\fill[stA] (1,-4) rectangle (2,-5);
\fill[stB] (2,-4) rectangle (5,-5);
\fill[stA] (1,-5) rectangle (2,-6);
\fill[stB] (2,-5) rectangle (3,-6);
\fill[stB] (5,-5) rectangle (6,-6);
\fill[stA] (1,-6) rectangle (2,-7);
\fill[stB] (2,-6) rectangle (7,-7);
\fill[stA] (1,-7) rectangle (2,-8);
\fill[stB] (2,-7) rectangle (3,-8);
\fill[stB] (7,-7) rectangle (8,-8);
\fill[stA] (1,-8) rectangle (2,-9);
\fill[stB] (2,-8) rectangle (4,-9);
\fill[stB] (6,-8) rectangle (9,-9);
\fill[stA] (1,-9) rectangle (2,-10);
\fill[stB] (2,-9) rectangle (3,-10);
\fill[stB] (4,-9) rectangle (6,-10);
\fill[stB] (9,-9) rectangle (10,-10);
\fill[stA] (1,-10) rectangle (2,-11);
\fill[stB] (2,-10) rectangle (3,-11);
\fill[stB] (6,-10) rectangle (7,-11);
\fill[stB] (8,-10) rectangle (11,-11);
\fill[stA] (1,-11) rectangle (2,-12);
\fill[stB] (2,-11) rectangle (4,-12);
\fill[stB] (5,-11) rectangle (7,-12);
\fill[stB] (11,-11) rectangle (12,-12);
\fill[stA] (1,-12) rectangle (2,-13);
\fill[stB] (2,-12) rectangle (3,-13);
\fill[stB] (7,-12) rectangle (8,-13);
\fill[stB] (10,-12) rectangle (13,-13);
\fill[stA] (1,-13) rectangle (2,-14);
\fill[stB] (2,-13) rectangle (4,-14);
\fill[stB] (6,-13) rectangle (10,-14);
\fill[stA] (13,-13) rectangle (14,-14);
\fill[stA] (1,-14) rectangle (2,-15);
\fill[stB] (2,-14) rectangle (3,-15);
\fill[stB] (4,-14) rectangle (6,-15);
\fill[stB] (10,-14) rectangle (11,-15);
\fill[stG] (12,-14) rectangle (13,-15);
\fill[stA] (13,-14) rectangle (14,-15);
\fill[stA] (1,-15) rectangle (2,-16);
\fill[stB] (2,-15) rectangle (3,-16);
\fill[stB] (6,-15) rectangle (7,-16);
\fill[stB] (9,-15) rectangle (11,-16);
\fill[stA] (11,-15) rectangle (12,-16);
\fill[stB] (12,-15) rectangle (14,-16);
\fill[stA] (1,-16) rectangle (2,-17);
\fill[stB] (2,-16) rectangle (4,-17);
\fill[stB] (5,-16) rectangle (9,-17);
\fill[stA] (10,-16) rectangle (12,-17);
\fill[stB] (12,-16) rectangle (13,-17);
\fill[stA] (13,-16) rectangle (14,-17);
\fill[stA] (1,-17) rectangle (2,-18);
\fill[stB] (2,-17) rectangle (3,-18);
\fill[stA] (9,-17) rectangle (10,-18);
\fill[stB] (10,-17) rectangle (11,-18);
\fill[stA] (11,-17) rectangle (12,-18);
\fill[stG] (12,-17) rectangle (13,-18);
\fill[stA] (1,-18) rectangle (2,-19);
\fill[stB] (2,-18) rectangle (4,-19);
\fill[stG] (8,-18) rectangle (9,-19);
\fill[stA] (9,-18) rectangle (10,-19);
\fill[stG] (10,-18) rectangle (11,-19);
\fill[stA] (11,-18) rectangle (13,-19);
\fill[stG] (13,-18) rectangle (14,-19);
\fill[stA] (1,-19) rectangle (2,-20);
\fill[stB] (2,-19) rectangle (3,-20);
\fill[stB] (4,-19) rectangle (5,-20);
\fill[stG] (7,-19) rectangle (8,-20);
\fill[stB] (8,-19) rectangle (9,-20);
\fill[stG] (9,-19) rectangle (11,-20);
\fill[stA] (12,-19) rectangle (13,-20);
\fill[stB] (13,-19) rectangle (14,-20);
\fill[stA] (1,-20) rectangle (2,-21);
\fill[stB] (2,-20) rectangle (3,-21);
\fill[stB] (4,-20) rectangle (6,-21);
\fill[stG] (6,-20) rectangle (7,-21);
\fill[stA] (7,-20) rectangle (9,-21);
\fill[stG] (9,-20) rectangle (10,-21);
\fill[stB] (10,-20) rectangle (11,-21);
\fill[stA] (11,-20) rectangle (13,-21);
\fill[stG] (13,-20) rectangle (14,-21);
\fill[stA] (1,-21) rectangle (2,-22);
\fill[stB] (2,-21) rectangle (3,-22);
\fill[stA] (5,-21) rectangle (7,-22);
\fill[stA] (8,-21) rectangle (10,-22);
\fill[stB] (10,-21) rectangle (11,-22);
\fill[stG] (11,-21) rectangle (12,-22);
\fill[stA] (12,-21) rectangle (13,-22);
\fill[stB] (13,-21) rectangle (14,-22);
\fill[stA] (1,-22) rectangle (2,-23);
\fill[stB] (2,-22) rectangle (4,-23);
\fill[stG] (4,-22) rectangle (5,-23);
\fill[stB] (5,-22) rectangle (6,-23);
\fill[stG] (7,-22) rectangle (8,-23);
\fill[stB] (8,-22) rectangle (9,-23);
\fill[stA] (9,-22) rectangle (10,-23);
\fill[stB] (10,-22) rectangle (11,-23);
\fill[stA] (11,-22) rectangle (13,-23);
\fill[stG] (13,-22) rectangle (14,-23);
\fill[stA] (1,-23) rectangle (2,-24);
\fill[stB] (2,-23) rectangle (3,-24);
\fill[stA] (3,-23) rectangle (4,-24);
\fill[stB] (5,-23) rectangle (6,-24);
\fill[stA] (6,-23) rectangle (8,-24);
\fill[stB] (8,-23) rectangle (9,-24);
\fill[stA] (9,-23) rectangle (10,-24);
\fill[stG] (10,-23) rectangle (12,-24);
\fill[stA] (12,-23) rectangle (13,-24);
\fill[stB] (13,-23) rectangle (14,-24);
\fill[stA] (1,-24) rectangle (2,-25);
\fill[stG] (2,-24) rectangle (3,-25);
\fill[stA] (3,-24) rectangle (4,-25);
\fill[stG] (4,-24) rectangle (5,-25);
\fill[stA] (5,-24) rectangle (6,-25);
\fill[stG] (6,-24) rectangle (7,-25);
\fill[stA] (7,-24) rectangle (8,-25);
\fill[stG] (8,-24) rectangle (9,-25);
\fill[stA] (9,-24) rectangle (10,-25);
\fill[stG] (10,-24) rectangle (12,-25);
\fill[stA] (12,-24) rectangle (13,-25);
\fill[stG] (13,-24) rectangle (14,-25);
\fill[stG] (1,-25) rectangle (13,-26);
\fill[stB] (13,-25) rectangle (14,-26);
\fill[stF] (0,-26) rectangle (14,-27);

    \draw[grid,line width=0.3pt] (0,0) rectangle (14,-27);
    \node[font=\small,anchor=south] at (7,0.6) {$n=14$};
    \foreach \y/\l in {0/0,13/13,26/26}
      \node[font=\scriptsize,anchor=east] at (-0.5,-\y-0.5) {$\l$};
    \foreach \x/\l in {1/1,14/14}
      \node[font=\scriptsize,anchor=north] at (\x-0.5,-27.3) {$\l$};
  \end{scope}
  % the line of length 15, up to its first firing cell (time 22)
  \begin{scope}[x=1.55mm,y=1.55mm,shift={(19,0)}]
    %% generated by figures.py; do not edit
\fill[stG] (0,0) rectangle (1,-1);
\fill[stA] (1,-1) rectangle (2,-2);
\fill[stB] (1,-2) rectangle (3,-3);
\fill[stG] (0,-3) rectangle (1,-4);
\fill[stB] (3,-3) rectangle (4,-4);
\fill[stA] (1,-4) rectangle (2,-5);
\fill[stB] (2,-4) rectangle (5,-5);
\fill[stA] (1,-5) rectangle (2,-6);
\fill[stB] (2,-5) rectangle (3,-6);
\fill[stB] (5,-5) rectangle (6,-6);
\fill[stA] (1,-6) rectangle (2,-7);
\fill[stB] (2,-6) rectangle (7,-7);
\fill[stA] (1,-7) rectangle (2,-8);
\fill[stB] (2,-7) rectangle (3,-8);
\fill[stB] (7,-7) rectangle (8,-8);
\fill[stA] (1,-8) rectangle (2,-9);
\fill[stB] (2,-8) rectangle (4,-9);
\fill[stB] (6,-8) rectangle (9,-9);
\fill[stA] (1,-9) rectangle (2,-10);
\fill[stB] (2,-9) rectangle (3,-10);
\fill[stB] (4,-9) rectangle (6,-10);
\fill[stB] (9,-9) rectangle (10,-10);
\fill[stA] (1,-10) rectangle (2,-11);
\fill[stB] (2,-10) rectangle (3,-11);
\fill[stB] (6,-10) rectangle (7,-11);
\fill[stB] (8,-10) rectangle (11,-11);
\fill[stA] (1,-11) rectangle (2,-12);
\fill[stB] (2,-11) rectangle (4,-12);
\fill[stB] (5,-11) rectangle (7,-12);
\fill[stB] (11,-11) rectangle (12,-12);
\fill[stA] (1,-12) rectangle (2,-13);
\fill[stB] (2,-12) rectangle (3,-13);
\fill[stB] (7,-12) rectangle (8,-13);
\fill[stB] (10,-12) rectangle (13,-13);
\fill[stA] (1,-13) rectangle (2,-14);
\fill[stB] (2,-13) rectangle (4,-14);
\fill[stB] (6,-13) rectangle (10,-14);
\fill[stB] (13,-13) rectangle (14,-14);
\fill[stA] (1,-14) rectangle (2,-15);
\fill[stB] (2,-14) rectangle (3,-15);
\fill[stB] (4,-14) rectangle (6,-15);
\fill[stB] (10,-14) rectangle (11,-15);
\fill[stB] (12,-14) rectangle (14,-15);
\fill[stA] (14,-14) rectangle (15,-15);
\fill[stA] (1,-15) rectangle (2,-16);
\fill[stB] (2,-15) rectangle (3,-16);
\fill[stB] (6,-15) rectangle (7,-16);
\fill[stB] (9,-15) rectangle (11,-16);
\fill[stA] (13,-15) rectangle (14,-16);
\fill[stA] (1,-16) rectangle (2,-17);
\fill[stB] (2,-16) rectangle (4,-17);
\fill[stB] (5,-16) rectangle (9,-17);
\fill[stB] (11,-16) rectangle (12,-17);
\fill[stG] (12,-16) rectangle (13,-17);
\fill[stB] (13,-16) rectangle (14,-17);
\fill[stG] (14,-16) rectangle (15,-17);
\fill[stA] (1,-17) rectangle (2,-18);
\fill[stB] (2,-17) rectangle (3,-18);
\fill[stB] (9,-17) rectangle (11,-18);
\fill[stG] (11,-17) rectangle (12,-18);
\fill[stA] (13,-17) rectangle (14,-18);
\fill[stB] (14,-17) rectangle (15,-18);
\fill[stA] (1,-18) rectangle (2,-19);
\fill[stB] (2,-18) rectangle (4,-19);
\fill[stB] (8,-18) rectangle (9,-19);
\fill[stA] (10,-18) rectangle (14,-19);
\fill[stG] (14,-18) rectangle (15,-19);
\fill[stA] (1,-19) rectangle (2,-20);
\fill[stB] (2,-19) rectangle (3,-20);
\fill[stB] (4,-19) rectangle (5,-20);
\fill[stB] (7,-19) rectangle (9,-20);
\fill[stA] (9,-19) rectangle (10,-20);
\fill[stB] (10,-19) rectangle (11,-20);
\fill[stG] (11,-19) rectangle (13,-20);
\fill[stA] (13,-19) rectangle (14,-20);
\fill[stB] (14,-19) rectangle (15,-20);
\fill[stA] (1,-20) rectangle (2,-21);
\fill[stB] (2,-20) rectangle (3,-21);
\fill[stB] (4,-20) rectangle (7,-21);
\fill[stA] (8,-20) rectangle (10,-21);
\fill[stB] (10,-20) rectangle (11,-21);
\fill[stG] (11,-20) rectangle (13,-21);
\fill[stA] (13,-20) rectangle (14,-21);
\fill[stG] (14,-20) rectangle (15,-21);
\fill[stA] (1,-21) rectangle (2,-22);
\fill[stB] (2,-21) rectangle (3,-22);
\fill[stA] (7,-21) rectangle (8,-22);
\fill[stB] (8,-21) rectangle (9,-22);
\fill[stA] (9,-21) rectangle (10,-22);
\fill[stB] (10,-21) rectangle (11,-22);
\fill[stG] (11,-21) rectangle (14,-22);
\fill[stB] (14,-21) rectangle (15,-22);
\fill[stA] (1,-22) rectangle (2,-23);
\fill[stB] (2,-22) rectangle (4,-23);
\fill[stG] (6,-22) rectangle (7,-23);
\fill[stA] (7,-22) rectangle (8,-23);
\fill[stG] (8,-22) rectangle (9,-23);
\fill[stA] (9,-22) rectangle (10,-23);
\fill[stB] (10,-22) rectangle (11,-23);
\fill[stG] (11,-22) rectangle (12,-23);
\fill[stF] (12,-22) rectangle (15,-23);

    \draw[grid,line width=0.3pt] (0,0) rectangle (15,-23);
    \draw[ink,line width=1.0pt] (12,-22) rectangle (15,-23);
    \node[font=\small,anchor=south] at (7.5,0.6) {$n=15$};
    \node[font=\scriptsize,anchor=north,align=center] at (7.5,-23.4) {cells $13$--$15$ fire\\ at $t=22$};
    \foreach \y/\l in {0/0,14/14,22/22}
      \node[font=\scriptsize,anchor=west] at (15.4,-\y-0.5) {$\l$};
  \end{scope}
  % the half-line, 0 <= t <= 63
  \begin{scope}[x=0.92mm,y=0.92mm,shift={(70,0)}]
    %% generated by figures.py; do not edit
\fill[stG] (0,0) rectangle (1,-1);
\fill[stA] (1,-1) rectangle (2,-2);
\fill[stB] (1,-2) rectangle (3,-3);
\fill[stG] (0,-3) rectangle (1,-4);
\fill[stB] (3,-3) rectangle (4,-4);
\fill[stA] (1,-4) rectangle (2,-5);
\fill[stB] (2,-4) rectangle (5,-5);
\fill[stA] (1,-5) rectangle (2,-6);
\fill[stB] (2,-5) rectangle (3,-6);
\fill[stB] (5,-5) rectangle (6,-6);
\fill[stA] (1,-6) rectangle (2,-7);
\fill[stB] (2,-6) rectangle (7,-7);
\fill[stA] (1,-7) rectangle (2,-8);
\fill[stB] (2,-7) rectangle (3,-8);
\fill[stB] (7,-7) rectangle (8,-8);
\fill[stA] (1,-8) rectangle (2,-9);
\fill[stB] (2,-8) rectangle (4,-9);
\fill[stB] (6,-8) rectangle (9,-9);
\fill[stA] (1,-9) rectangle (2,-10);
\fill[stB] (2,-9) rectangle (3,-10);
\fill[stB] (4,-9) rectangle (6,-10);
\fill[stB] (9,-9) rectangle (10,-10);
\fill[stA] (1,-10) rectangle (2,-11);
\fill[stB] (2,-10) rectangle (3,-11);
\fill[stB] (6,-10) rectangle (7,-11);
\fill[stB] (8,-10) rectangle (11,-11);
\fill[stA] (1,-11) rectangle (2,-12);
\fill[stB] (2,-11) rectangle (4,-12);
\fill[stB] (5,-11) rectangle (7,-12);
\fill[stB] (11,-11) rectangle (12,-12);
\fill[stA] (1,-12) rectangle (2,-13);
\fill[stB] (2,-12) rectangle (3,-13);
\fill[stB] (7,-12) rectangle (8,-13);
\fill[stB] (10,-12) rectangle (13,-13);
\fill[stA] (1,-13) rectangle (2,-14);
\fill[stB] (2,-13) rectangle (4,-14);
\fill[stB] (6,-13) rectangle (10,-14);
\fill[stB] (13,-13) rectangle (14,-14);
\fill[stA] (1,-14) rectangle (2,-15);
\fill[stB] (2,-14) rectangle (3,-15);
\fill[stB] (4,-14) rectangle (6,-15);
\fill[stB] (10,-14) rectangle (11,-15);
\fill[stB] (12,-14) rectangle (15,-15);
\fill[stA] (1,-15) rectangle (2,-16);
\fill[stB] (2,-15) rectangle (3,-16);
\fill[stB] (6,-15) rectangle (7,-16);
\fill[stB] (9,-15) rectangle (11,-16);
\fill[stB] (15,-15) rectangle (16,-16);
\fill[stA] (1,-16) rectangle (2,-17);
\fill[stB] (2,-16) rectangle (4,-17);
\fill[stB] (5,-16) rectangle (9,-17);
\fill[stB] (11,-16) rectangle (12,-17);
\fill[stB] (14,-16) rectangle (17,-17);
\fill[stA] (1,-17) rectangle (2,-18);
\fill[stB] (2,-17) rectangle (3,-18);
\fill[stB] (9,-17) rectangle (14,-18);
\fill[stB] (17,-17) rectangle (18,-18);
\fill[stA] (1,-18) rectangle (2,-19);
\fill[stB] (2,-18) rectangle (4,-19);
\fill[stB] (8,-18) rectangle (9,-19);
\fill[stB] (14,-18) rectangle (15,-19);
\fill[stB] (16,-18) rectangle (19,-19);
\fill[stA] (1,-19) rectangle (2,-20);
\fill[stB] (2,-19) rectangle (3,-20);
\fill[stB] (4,-19) rectangle (5,-20);
\fill[stB] (7,-19) rectangle (10,-20);
\fill[stB] (13,-19) rectangle (15,-20);
\fill[stB] (19,-19) rectangle (20,-20);
\fill[stA] (1,-20) rectangle (2,-21);
\fill[stB] (2,-20) rectangle (3,-21);
\fill[stB] (4,-20) rectangle (7,-21);
\fill[stB] (10,-20) rectangle (11,-21);
\fill[stB] (12,-20) rectangle (13,-21);
\fill[stB] (15,-20) rectangle (16,-21);
\fill[stB] (18,-20) rectangle (21,-21);
\fill[stA] (1,-21) rectangle (2,-22);
\fill[stB] (2,-21) rectangle (3,-22);
\fill[stB] (7,-21) rectangle (8,-22);
\fill[stB] (9,-21) rectangle (11,-22);
\fill[stB] (12,-21) rectangle (18,-22);
\fill[stB] (21,-21) rectangle (22,-22);
\fill[stA] (1,-22) rectangle (2,-23);
\fill[stB] (2,-22) rectangle (4,-23);
\fill[stB] (6,-22) rectangle (8,-23);
\fill[stB] (18,-22) rectangle (19,-23);
\fill[stB] (20,-22) rectangle (23,-23);
\fill[stA] (1,-23) rectangle (2,-24);
\fill[stB] (2,-23) rectangle (3,-24);
\fill[stB] (4,-23) rectangle (6,-24);
\fill[stB] (8,-23) rectangle (9,-24);
\fill[stB] (17,-23) rectangle (19,-24);
\fill[stB] (23,-23) rectangle (24,-24);
\fill[stA] (1,-24) rectangle (2,-25);
\fill[stB] (2,-24) rectangle (3,-25);
\fill[stB] (6,-24) rectangle (10,-25);
\fill[stB] (16,-24) rectangle (17,-25);
\fill[stB] (19,-24) rectangle (20,-25);
\fill[stB] (22,-24) rectangle (25,-25);
\fill[stA] (1,-25) rectangle (2,-26);
\fill[stB] (2,-25) rectangle (4,-26);
\fill[stB] (5,-25) rectangle (6,-26);
\fill[stB] (10,-25) rectangle (11,-26);
\fill[stB] (15,-25) rectangle (22,-26);
\fill[stB] (25,-25) rectangle (26,-26);
\fill[stA] (1,-26) rectangle (2,-27);
\fill[stB] (2,-26) rectangle (3,-27);
\fill[stB] (5,-26) rectangle (7,-27);
\fill[stB] (9,-26) rectangle (12,-27);
\fill[stB] (14,-26) rectangle (15,-27);
\fill[stB] (22,-26) rectangle (23,-27);
\fill[stB] (24,-26) rectangle (27,-27);
\fill[stA] (1,-27) rectangle (2,-28);
\fill[stB] (2,-27) rectangle (5,-28);
\fill[stB] (7,-27) rectangle (9,-28);
\fill[stB] (12,-27) rectangle (16,-28);
\fill[stB] (21,-27) rectangle (23,-28);
\fill[stB] (27,-27) rectangle (28,-28);
\fill[stA] (1,-28) rectangle (2,-29);
\fill[stB] (2,-28) rectangle (3,-29);
\fill[stB] (5,-28) rectangle (7,-29);
\fill[stB] (9,-28) rectangle (10,-29);
\fill[stB] (11,-28) rectangle (12,-29);
\fill[stB] (16,-28) rectangle (17,-29);
\fill[stB] (20,-28) rectangle (21,-29);
\fill[stB] (23,-28) rectangle (24,-29);
\fill[stB] (26,-28) rectangle (29,-29);
\fill[stA] (1,-29) rectangle (2,-30);
\fill[stB] (2,-29) rectangle (5,-30);
\fill[stB] (7,-29) rectangle (10,-30);
\fill[stB] (11,-29) rectangle (13,-30);
\fill[stB] (15,-29) rectangle (18,-30);
\fill[stB] (19,-29) rectangle (26,-30);
\fill[stB] (29,-29) rectangle (30,-30);
\fill[stA] (1,-30) rectangle (2,-31);
\fill[stB] (2,-30) rectangle (3,-31);
\fill[stB] (5,-30) rectangle (7,-31);
\fill[stB] (13,-30) rectangle (15,-31);
\fill[stB] (26,-30) rectangle (27,-31);
\fill[stB] (28,-30) rectangle (31,-31);
\fill[stA] (1,-31) rectangle (2,-32);
\fill[stB] (2,-31) rectangle (5,-32);
\fill[stB] (7,-31) rectangle (8,-32);
\fill[stB] (12,-31) rectangle (13,-32);
\fill[stB] (15,-31) rectangle (16,-32);
\fill[stB] (25,-31) rectangle (27,-32);
\fill[stB] (31,-31) rectangle (32,-32);
\fill[stA] (1,-32) rectangle (2,-33);
\fill[stB] (2,-32) rectangle (3,-33);
\fill[stB] (5,-32) rectangle (9,-33);
\fill[stB] (11,-32) rectangle (17,-33);
\fill[stB] (24,-32) rectangle (25,-33);
\fill[stB] (27,-32) rectangle (28,-33);
\fill[stB] (30,-32) rectangle (33,-33);
\fill[stA] (1,-33) rectangle (2,-34);
\fill[stB] (2,-33) rectangle (5,-34);
\fill[stB] (9,-33) rectangle (11,-34);
\fill[stB] (17,-33) rectangle (18,-34);
\fill[stB] (23,-33) rectangle (30,-34);
\fill[stB] (33,-33) rectangle (34,-34);
\fill[stA] (1,-34) rectangle (2,-35);
\fill[stB] (2,-34) rectangle (3,-35);
\fill[stB] (5,-34) rectangle (6,-35);
\fill[stB] (8,-34) rectangle (9,-35);
\fill[stB] (11,-34) rectangle (12,-35);
\fill[stB] (16,-34) rectangle (19,-35);
\fill[stB] (22,-34) rectangle (23,-35);
\fill[stB] (30,-34) rectangle (31,-35);
\fill[stB] (32,-34) rectangle (35,-35);
\fill[stA] (1,-35) rectangle (2,-36);
\fill[stB] (2,-35) rectangle (13,-36);
\fill[stB] (15,-35) rectangle (16,-36);
\fill[stB] (19,-35) rectangle (20,-36);
\fill[stB] (21,-35) rectangle (24,-36);
\fill[stB] (29,-35) rectangle (31,-36);
\fill[stB] (35,-35) rectangle (36,-36);
\fill[stA] (1,-36) rectangle (2,-37);
\fill[stB] (2,-36) rectangle (3,-37);
\fill[stB] (13,-36) rectangle (17,-37);
\fill[stB] (18,-36) rectangle (20,-37);
\fill[stB] (24,-36) rectangle (25,-37);
\fill[stB] (28,-36) rectangle (29,-37);
\fill[stB] (31,-36) rectangle (32,-37);
\fill[stB] (34,-36) rectangle (37,-37);
\fill[stA] (1,-37) rectangle (2,-38);
\fill[stB] (2,-37) rectangle (4,-38);
\fill[stB] (12,-37) rectangle (13,-38);
\fill[stB] (20,-37) rectangle (21,-38);
\fill[stB] (23,-37) rectangle (26,-38);
\fill[stB] (27,-37) rectangle (34,-38);
\fill[stB] (37,-37) rectangle (38,-38);
\fill[stA] (1,-38) rectangle (2,-39);
\fill[stB] (2,-38) rectangle (3,-39);
\fill[stB] (4,-38) rectangle (5,-39);
\fill[stB] (11,-38) rectangle (14,-39);
\fill[stB] (19,-38) rectangle (23,-39);
\fill[stB] (34,-38) rectangle (35,-39);
\fill[stB] (36,-38) rectangle (39,-39);
\fill[stA] (1,-39) rectangle (2,-40);
\fill[stB] (2,-39) rectangle (3,-40);
\fill[stB] (4,-39) rectangle (6,-40);
\fill[stB] (10,-39) rectangle (11,-40);
\fill[stB] (14,-39) rectangle (15,-40);
\fill[stB] (18,-39) rectangle (19,-40);
\fill[stB] (23,-39) rectangle (24,-40);
\fill[stB] (33,-39) rectangle (35,-40);
\fill[stB] (39,-39) rectangle (40,-40);
\fill[stA] (1,-40) rectangle (2,-41);
\fill[stB] (2,-40) rectangle (3,-41);
\fill[stB] (6,-40) rectangle (7,-41);
\fill[stB] (9,-40) rectangle (12,-41);
\fill[stB] (13,-40) rectangle (16,-41);
\fill[stB] (17,-40) rectangle (20,-41);
\fill[stB] (22,-40) rectangle (25,-41);
\fill[stB] (32,-40) rectangle (33,-41);
\fill[stB] (35,-40) rectangle (36,-41);
\fill[stB] (38,-40) rectangle (41,-41);
\fill[stA] (1,-41) rectangle (2,-42);
\fill[stB] (2,-41) rectangle (4,-42);
\fill[stB] (5,-41) rectangle (9,-42);
\fill[stB] (20,-41) rectangle (22,-42);
\fill[stB] (25,-41) rectangle (26,-42);
\fill[stB] (31,-41) rectangle (38,-42);
\fill[stB] (41,-41) rectangle (42,-42);
\fill[stA] (1,-42) rectangle (2,-43);
\fill[stB] (2,-42) rectangle (3,-43);
\fill[stB] (9,-42) rectangle (10,-43);
\fill[stB] (19,-42) rectangle (20,-43);
\fill[stB] (22,-42) rectangle (23,-43);
\fill[stB] (24,-42) rectangle (27,-43);
\fill[stB] (30,-42) rectangle (31,-43);
\fill[stB] (38,-42) rectangle (39,-43);
\fill[stB] (40,-42) rectangle (43,-43);
\fill[stA] (1,-43) rectangle (2,-44);
\fill[stB] (2,-43) rectangle (4,-44);
\fill[stB] (8,-43) rectangle (11,-44);
\fill[stB] (18,-43) rectangle (23,-44);
\fill[stB] (27,-43) rectangle (28,-44);
\fill[stB] (29,-43) rectangle (32,-44);
\fill[stB] (37,-43) rectangle (39,-44);
\fill[stB] (43,-43) rectangle (44,-44);
\fill[stA] (1,-44) rectangle (2,-45);
\fill[stB] (2,-44) rectangle (3,-45);
\fill[stB] (4,-44) rectangle (5,-45);
\fill[stB] (7,-44) rectangle (8,-45);
\fill[stB] (11,-44) rectangle (12,-45);
\fill[stB] (17,-44) rectangle (18,-45);
\fill[stB] (23,-44) rectangle (24,-45);
\fill[stB] (26,-44) rectangle (28,-45);
\fill[stB] (32,-44) rectangle (33,-45);
\fill[stB] (36,-44) rectangle (37,-45);
\fill[stB] (39,-44) rectangle (40,-45);
\fill[stB] (42,-44) rectangle (45,-45);
\fill[stA] (1,-45) rectangle (2,-46);
\fill[stB] (2,-45) rectangle (3,-46);
\fill[stB] (4,-45) rectangle (9,-46);
\fill[stB] (10,-45) rectangle (13,-46);
\fill[stB] (16,-45) rectangle (19,-46);
\fill[stB] (22,-45) rectangle (26,-46);
\fill[stB] (28,-45) rectangle (29,-46);
\fill[stB] (31,-45) rectangle (34,-46);
\fill[stB] (35,-45) rectangle (42,-46);
\fill[stB] (45,-45) rectangle (46,-46);
\fill[stA] (1,-46) rectangle (2,-47);
\fill[stB] (2,-46) rectangle (3,-47);
\fill[stB] (13,-46) rectangle (14,-47);
\fill[stB] (15,-46) rectangle (16,-47);
\fill[stB] (19,-46) rectangle (20,-47);
\fill[stB] (21,-46) rectangle (22,-47);
\fill[stB] (26,-46) rectangle (31,-47);
\fill[stB] (42,-46) rectangle (43,-47);
\fill[stB] (44,-46) rectangle (47,-47);
\fill[stA] (1,-47) rectangle (2,-48);
\fill[stB] (2,-47) rectangle (4,-48);
\fill[stB] (12,-47) rectangle (14,-48);
\fill[stB] (15,-47) rectangle (17,-48);
\fill[stB] (18,-47) rectangle (20,-48);
\fill[stB] (21,-47) rectangle (23,-48);
\fill[stB] (25,-47) rectangle (26,-48);
\fill[stB] (31,-47) rectangle (32,-48);
\fill[stB] (41,-47) rectangle (43,-48);
\fill[stB] (47,-47) rectangle (48,-48);
\fill[stA] (1,-48) rectangle (2,-49);
\fill[stB] (2,-48) rectangle (3,-49);
\fill[stB] (4,-48) rectangle (5,-49);
\fill[stB] (11,-48) rectangle (12,-49);
\fill[stB] (23,-48) rectangle (27,-49);
\fill[stB] (30,-48) rectangle (33,-49);
\fill[stB] (40,-48) rectangle (41,-49);
\fill[stB] (43,-48) rectangle (44,-49);
\fill[stB] (46,-48) rectangle (49,-49);
\fill[stA] (1,-49) rectangle (2,-50);
\fill[stB] (2,-49) rectangle (3,-50);
\fill[stB] (4,-49) rectangle (6,-50);
\fill[stB] (10,-49) rectangle (13,-50);
\fill[stB] (22,-49) rectangle (23,-50);
\fill[stB] (27,-49) rectangle (28,-50);
\fill[stB] (29,-49) rectangle (30,-50);
\fill[stB] (33,-49) rectangle (34,-50);
\fill[stB] (39,-49) rectangle (46,-50);
\fill[stB] (49,-49) rectangle (50,-50);
\fill[stA] (1,-50) rectangle (2,-51);
\fill[stB] (2,-50) rectangle (3,-51);
\fill[stB] (6,-50) rectangle (7,-51);
\fill[stB] (9,-50) rectangle (10,-51);
\fill[stB] (13,-50) rectangle (14,-51);
\fill[stB] (21,-50) rectangle (24,-51);
\fill[stB] (26,-50) rectangle (28,-51);
\fill[stB] (29,-50) rectangle (31,-51);
\fill[stB] (32,-50) rectangle (35,-51);
\fill[stB] (38,-50) rectangle (39,-51);
\fill[stB] (46,-50) rectangle (47,-51);
\fill[stB] (48,-50) rectangle (51,-51);
\fill[stA] (1,-51) rectangle (2,-52);
\fill[stB] (2,-51) rectangle (4,-52);
\fill[stB] (5,-51) rectangle (11,-52);
\fill[stB] (12,-51) rectangle (15,-52);
\fill[stB] (20,-51) rectangle (21,-52);
\fill[stB] (24,-51) rectangle (26,-52);
\fill[stB] (35,-51) rectangle (36,-52);
\fill[stB] (37,-51) rectangle (40,-52);
\fill[stB] (45,-51) rectangle (47,-52);
\fill[stB] (51,-51) rectangle (52,-52);
\fill[stA] (1,-52) rectangle (2,-53);
\fill[stB] (2,-52) rectangle (3,-53);
\fill[stB] (15,-52) rectangle (16,-53);
\fill[stB] (19,-52) rectangle (22,-53);
\fill[stB] (23,-52) rectangle (24,-53);
\fill[stB] (26,-52) rectangle (27,-53);
\fill[stB] (34,-52) rectangle (36,-53);
\fill[stB] (40,-52) rectangle (41,-53);
\fill[stB] (44,-52) rectangle (45,-53);
\fill[stB] (47,-52) rectangle (48,-53);
\fill[stB] (50,-52) rectangle (53,-53);
\fill[stA] (1,-53) rectangle (2,-54);
\fill[stB] (2,-53) rectangle (4,-54);
\fill[stB] (14,-53) rectangle (17,-54);
\fill[stB] (18,-53) rectangle (19,-54);
\fill[stB] (23,-53) rectangle (28,-54);
\fill[stB] (33,-53) rectangle (34,-54);
\fill[stB] (36,-53) rectangle (37,-54);
\fill[stB] (39,-53) rectangle (42,-54);
\fill[stB] (43,-53) rectangle (50,-54);
\fill[stB] (53,-53) rectangle (54,-54);
\fill[stA] (1,-54) rectangle (2,-55);
\fill[stB] (2,-54) rectangle (3,-55);
\fill[stB] (4,-54) rectangle (5,-55);
\fill[stB] (13,-54) rectangle (14,-55);
\fill[stB] (18,-54) rectangle (20,-55);
\fill[stB] (22,-54) rectangle (23,-55);
\fill[stB] (28,-54) rectangle (29,-55);
\fill[stB] (32,-54) rectangle (39,-55);
\fill[stB] (50,-54) rectangle (51,-55);
\fill[stB] (52,-54) rectangle (55,-55);
\fill[stA] (1,-55) rectangle (2,-56);
\fill[stB] (2,-55) rectangle (3,-56);
\fill[stB] (4,-55) rectangle (6,-56);
\fill[stB] (12,-55) rectangle (15,-56);
\fill[stB] (17,-55) rectangle (18,-56);
\fill[stB] (20,-55) rectangle (24,-56);
\fill[stB] (27,-55) rectangle (30,-56);
\fill[stB] (31,-55) rectangle (32,-56);
\fill[stB] (39,-55) rectangle (40,-56);
\fill[stB] (49,-55) rectangle (51,-56);
\fill[stB] (55,-55) rectangle (56,-56);
\fill[stA] (1,-56) rectangle (2,-57);
\fill[stB] (2,-56) rectangle (3,-57);
\fill[stB] (6,-56) rectangle (7,-57);
\fill[stB] (11,-56) rectangle (12,-57);
\fill[stB] (15,-56) rectangle (20,-57);
\fill[stB] (24,-56) rectangle (25,-57);
\fill[stB] (26,-56) rectangle (27,-57);
\fill[stB] (31,-56) rectangle (33,-57);
\fill[stB] (38,-56) rectangle (41,-57);
\fill[stB] (48,-56) rectangle (49,-57);
\fill[stB] (51,-56) rectangle (52,-57);
\fill[stB] (54,-56) rectangle (57,-57);
\fill[stA] (1,-57) rectangle (2,-58);
\fill[stB] (2,-57) rectangle (4,-58);
\fill[stB] (5,-57) rectangle (8,-58);
\fill[stB] (10,-57) rectangle (13,-58);
\fill[stB] (14,-57) rectangle (15,-58);
\fill[stB] (20,-57) rectangle (21,-58);
\fill[stB] (23,-57) rectangle (25,-58);
\fill[stB] (26,-57) rectangle (28,-58);
\fill[stB] (30,-57) rectangle (31,-58);
\fill[stB] (33,-57) rectangle (34,-58);
\fill[stB] (37,-57) rectangle (38,-58);
\fill[stB] (41,-57) rectangle (42,-58);
\fill[stB] (47,-57) rectangle (54,-58);
\fill[stB] (57,-57) rectangle (58,-58);
\fill[stA] (1,-58) rectangle (2,-59);
\fill[stB] (2,-58) rectangle (3,-59);
\fill[stB] (8,-58) rectangle (10,-59);
\fill[stB] (14,-58) rectangle (16,-59);
\fill[stB] (19,-58) rectangle (23,-59);
\fill[stB] (28,-58) rectangle (35,-59);
\fill[stB] (36,-58) rectangle (39,-59);
\fill[stB] (40,-58) rectangle (43,-59);
\fill[stB] (46,-58) rectangle (47,-59);
\fill[stB] (54,-58) rectangle (55,-59);
\fill[stB] (56,-58) rectangle (59,-59);
\fill[stA] (1,-59) rectangle (2,-60);
\fill[stB] (2,-59) rectangle (4,-60);
\fill[stB] (7,-59) rectangle (8,-60);
\fill[stB] (10,-59) rectangle (11,-60);
\fill[stB] (13,-59) rectangle (14,-60);
\fill[stB] (16,-59) rectangle (17,-60);
\fill[stB] (18,-59) rectangle (19,-60);
\fill[stB] (23,-59) rectangle (24,-60);
\fill[stB] (27,-59) rectangle (28,-60);
\fill[stB] (43,-59) rectangle (44,-60);
\fill[stB] (45,-59) rectangle (48,-60);
\fill[stB] (53,-59) rectangle (55,-60);
\fill[stB] (59,-59) rectangle (60,-60);
\fill[stA] (1,-60) rectangle (2,-61);
\fill[stB] (2,-60) rectangle (3,-61);
\fill[stB] (4,-60) rectangle (5,-61);
\fill[stB] (6,-60) rectangle (17,-61);
\fill[stB] (18,-60) rectangle (20,-61);
\fill[stB] (22,-60) rectangle (25,-61);
\fill[stB] (26,-60) rectangle (29,-61);
\fill[stB] (42,-60) rectangle (44,-61);
\fill[stB] (48,-60) rectangle (49,-61);
\fill[stB] (52,-60) rectangle (53,-61);
\fill[stB] (55,-60) rectangle (56,-61);
\fill[stB] (58,-60) rectangle (61,-61);
\fill[stA] (1,-61) rectangle (2,-62);
\fill[stB] (2,-61) rectangle (3,-62);
\fill[stB] (4,-61) rectangle (5,-62);
\fill[stB] (20,-61) rectangle (22,-62);
\fill[stB] (29,-61) rectangle (30,-62);
\fill[stB] (41,-61) rectangle (42,-62);
\fill[stB] (44,-61) rectangle (45,-62);
\fill[stB] (47,-61) rectangle (50,-62);
\fill[stB] (51,-61) rectangle (58,-62);
\fill[stB] (61,-61) rectangle (62,-62);
\fill[stA] (1,-62) rectangle (2,-63);
\fill[stB] (2,-62) rectangle (3,-63);
\fill[stB] (4,-62) rectangle (6,-63);
\fill[stB] (19,-62) rectangle (20,-63);
\fill[stB] (22,-62) rectangle (23,-63);
\fill[stB] (28,-62) rectangle (31,-63);
\fill[stB] (40,-62) rectangle (47,-63);
\fill[stB] (58,-62) rectangle (59,-63);
\fill[stB] (60,-62) rectangle (63,-63);
\fill[stA] (1,-63) rectangle (2,-64);
\fill[stB] (2,-63) rectangle (3,-64);
\fill[stB] (6,-63) rectangle (7,-64);
\fill[stB] (18,-63) rectangle (24,-64);
\fill[stB] (27,-63) rectangle (28,-64);
\fill[stB] (31,-63) rectangle (32,-64);
\fill[stB] (39,-63) rectangle (40,-64);
\fill[stB] (47,-63) rectangle (48,-64);
\fill[stB] (57,-63) rectangle (59,-64);
\fill[stB] (63,-63) rectangle (64,-64);

    \draw[grid,line width=0.3pt] (0,0) rectangle (64,-64);
    \node[font=\small,anchor=south] at (32,1) {half-line};
    \foreach \y/\l in {0/0,63/63}
      \node[font=\scriptsize,anchor=east] at (-0.6,-\y-0.5) {$\l$};
  \end{scope}
  % legend
  \begin{scope}[shift={(0,-6.6)},font=\scriptsize]
    \foreach \c/\l [count=\k from 0] in {white/$\Lq$,stG/$\Gq$,stA/$\Aq$,stB/$\Bq$,stF/$\Fq$}
      \node[anchor=west,inner sep=0pt] at (\k*1.25,0) {\swatch{\c}\ \l};
  \end{scope}
\end{tikzpicture}

\caption{The rule $\delta_{14}$. Left: the line of length $14$ is synchronized
at time $26$. Middle: on the line of length $15$, the cells $13$, $14$ and $15$
fire at time $22$, that is, at the relative time $\tau=8$, six steps early.
Right: the half-line for $0\le t\le63$. From time $3$ on, the cells $i\ge3$
take only the values $\Lq$ and $\Bq$, and there $\delta_{14}$ acts as the
elementary cellular automaton $22$.}
\label{fig:delta14}
\end{figure}

\Cref{fig:delta14} shows how $\delta_{14}$ fails. For $3\le t<1500$ its
half-line takes the values $\Lq,\Gq$ at cell $1$, $\Lq,\Aq$ at cell $2$ and
$\Lq,\Bq$ at the cells $i\ge3$, and $\delta_{14}$ restricted to $\{\Lq,\Bq\}^3$
is the elementary rule $22$~\cite{Wolfram83}. In the simulated ranges, neither
barrier excludes this half-line: $\PUMP(n,n')$ holds for $4\le n<n'\le800$, and
\cref{lem:left} does not apply for $n\le5000$, $s\le3$ and $\ell\le24$. The
rules $\delta_{12}$ and $\delta_{13}$ synchronize the lines of lengths
$2,\dots,12$ and $2,\dots,13$, respectively, and no other line of length at most
$60$. Each of these seven rules has a line of some length $14\le n\le19$ with a
cell in state $\Fq$ at a relative time $\tau\le9$, so none of them satisfies
$\G_5(9,N_1)$ for $N_1\ge19$.

\subsection{Five States: Half-Lines of Low Complexity}\label{sec:germs}

The barriers constrain the half-line, so it is natural to ask how simple the
half-line of a five-state minimal-time solution can be. We measure this by the
number of distinct neighbourhoods that it uses in a fixed region. The
\emph{germ} below $M$ of a rule is its restriction to the neighbourhoods of the
cells $(t,i)$ of its half-line with $t\ge1$ and $t+i\le M$, and its
\emph{complexity} $c_M(\delta)$ is the number of these neighbourhoods other than
$(\Lq,\Lq,\Lq)$. We use $M=78$: the half-line below the anti-diagonal $t+i=78$
contains the input anti-diagonals of all lines of lengths $n\le40$, so the germ
below $78$ determines the conditions $\PUMP(n,n')$ for $n'\le40$. For example,
$c_{78}=57$ for Mazoyer's solution, $22$ for $\delta_{12}$ and $16$ for
$\delta_{14}$.

\begin{proposition}\label{prop:germs}
Every five-state minimal-time solution $\delta$ satisfies $c_{78}(\delta)\ge15$.
\end{proposition}

\begin{proof}
Suppose that $c_{78}(\delta)\le14$. Exchanging $\Aq$ and $\Bq$ if necessary, we
may assume that $\Aq$ occurs before $\Bq$ (if at all) in the order of
anti-diagonals and, within an anti-diagonal, of time. By \eqref{eq:neverfire}
and \cref{lem:return,lem:pumping}, the germ of $\delta$ below $78$ produces no
$\Fq$, never produces $\Lq$ at the front, and satisfies $\PUMP(n,n')$ for $4\le
n<n'\le40$. A depth-first search that chooses the value of each neighbourhood
at its first occurrence enumerates all such germs, visiting about $10^8$ nodes:
there are $244$, and an independent re-implementation finds the same $244$.
For $230$ of them, the formula $\MT_5(\{2,\dots,10\};78,39)$ with the germ fixed
is unsatisfiable. A single formula covers all $230$: it has a selector variable
for each germ, clauses that make each selector fix its germ, and a clause that
requires some selector; CaDiCaL refutes it with an LRAT proof that
\texttt{lrat-check} verifies. Each of the other $14$ germs determines
$C_\infty$ up to the anti-diagonal $3300$, and there \cref{lem:left} applies
with $s=1$: for some $n\le163$ and $\ell\le5$, the upper input anti-diagonal of
the line of length $n$ is $\ell$-periodic on $[1,Y]$ and the lower one on
$[1,Y-1]$ with $Y-1>16\ell$ (for example $n=28$, $\ell=1$ and $Y=19$), as an
independent simulation confirms. Hence no rule that contains one of the $244$
germs is a minimal-time solution.
\end{proof}

The bound is attained by rules that come close to minimal time: the germs below
$78$ of $\delta_{14}^{\mathrm c}$ and $\delta_{14}^{\mathrm d}$ have complexity
$15$, and both rules synchronize the lines of lengths $2,\dots,14$
(\cref{prop:nearmiss}).

The rules $\delta_{12}$ and $\delta_{13}$ have the same half-line on the cells
with $t+i\le78$, where they use the $22$ neighbourhoods
\begin{center}\small
\texttt{*GL>A GLL>G *AG>B AGL>L *BL>A BLG>L LGL>B *AL>G ALB>A LBG>L BGL>G}\\
\texttt{*GA>A GAL>L ALG>B LGG>G GGL>G BGG>L GGG>G ALL>A LLG>L ALA>A LAL>L}
\end{center}
besides $(\Lq,\Lq,\Lq)$; here \texttt{xyz>d} means $\delta(x,y,z)=d$. This
half-line consists of a front, a region of $\Gq$'s behind it and, beyond a
boundary of speed $1/3$, a region that alternates between $\Aq$ and $\Lq$,
periodically in time and space (\cref{fig:halflines}). By \cref{thm:left} no
minimal-time solution has such a half-line, and \cref{lem:left} gives an
explicit length.

\begin{corollary}\label{cor:germ12}
No five-state rule that has the $22$ transitions above synchronizes the line of
length $518$ in minimal time.
\end{corollary}

\begin{proof}
With these transitions and $\delta(\Lq,\Lq,\Lq)=\Lq$, no other neighbourhood
occurs on the cells with $t+i\le1034$, so every such rule has the same diagram
$C_\infty$ there. On the anti-diagonal $t+i=1034$ the cells $2,\dots,259$ are in
state $\Aq$, and on $t+i=1033$ the cells $2,\dots,258$ are in state $\Lq$.
\Cref{lem:left} with $k=5$, $n=518$, $s=2$, $\ell=1$ and $Y=259$ would give
$257\le4^4=256$.
\end{proof}

\subsection{Balzer's Conditions}\label{sec:balzer}

Balzer~\cite[p.~37]{Balzer67} reported that no five-state minimal-time
solution satisfies the following four conditions:
\begin{enumerate}[(B1)]
\item $\delta(x,\Gq,y)=\Gq$ unless $x,y\in\{\Gq,\ast\}$;
\item $\delta(I(z),I(y),I(x))=I(\delta(x,y,z))$ for an involution $I$ of the
states with $I(\ast)=\ast$;
\item $\delta(\Gq,\Gq,\Gq)=\delta(\ast,\Gq,\Gq)=\delta(\Gq,\Gq,\ast)=\Fq$, and no
other neighbourhood gives $\Fq$;
\item $\delta(\Gq,v,\Gq)=\Gq$ for all $v\in W\setminus\{\Gq\}$.
\end{enumerate}
In the \emph{weak reading}~\cite{Yunes93}, (B1) is required only for working
states $x,y$, and (B3) without its last clause. We take $I(\Fq)=\Fq$. Then
$I(\Gq)=\Gq$: otherwise $g:=I(\Gq)\in W\setminus\{\Gq\}$, and (B2) and (B4)
give $\delta(g,\Gq,g)=I(\delta(\Gq,g,\Gq))=I(\Gq)=g$, whereas (B1) gives
$\delta(g,\Gq,g)=\Gq$. So $I$ is one of the four involutions of
$\{\Lq,\Aq,\Bq\}$, and (B1)--(B4) become unit and binary clauses.

\begin{proposition}\label{prop:balzer}
No five-state rule that satisfies (B1)--(B4) in the weak reading, with
$I(\Fq)=\Fq$, synchronizes all lines of lengths $2,\dots,12$ in minimal time;
a fortiori, no rule in Balzer's form does. The lengths up to $11$ do not suffice
to prove this.
\end{proposition}

\begin{proof}
For each involution, the clauses of (B1)--(B4) are added to
$\MT_5(\{2,\dots,N\};2N-2)$. For $I=(\Aq\,\Bq)$ and $N=10$, and for $I=(\Lq\,\Aq)$
and $I=(\Lq\,\Bq)$ and $N=9$, CaDiCaL refutes the formula with an LRAT proof.
For $I$ the identity and $N=12$, Kissat refutes it; \textsf{drat-trim} turns its
DRAT proof into an LRAT proof. Each of the four proofs is verified by
\texttt{lrat-check}. The rule \texttt{results/balzer\_sym11.txt} satisfies
(B1)--(B4) in the strong reading with $I$ the identity and synchronizes every
line of length at most $11$ (direct simulation).
\end{proof}
